\chapter*{附录}
\addcontentsline{toc}{chapter}{附录}

\section{附录1 核心病虫害识别品类清单}

下表为项目 MVP 及迭代阶段计划覆盖的核心病虫害品类（示例，正式清单以知识库为准）：

\begin{table}[H]
  \centering
  \small
  \renewcommand{\arraystretch}{1.35}
  \caption{核心病虫害识别品类清单}
  \begin{tabular}{p{2.2cm}p{6.4cm}p{6.4cm}}
    \toprule
    作物 & 病害 & 虫害 \\
    \midrule
    水稻 & 稻瘟病、纹枯病、稻曲病、胡麻叶斑病、细菌性条斑病 & 稻纵卷叶螟、稻飞虱、二化螟 \\
    小麦 & 条锈病、叶锈病、赤霉病、白粉病 & 蚜虫、麦蜘蛛、吸浆虫 \\
    玉米 & 大斑病、小斑病、锈病、茎腐病 & 玉米螟、草地贪夜蛾、黏虫 \\
    番茄 & 早疫病、晚疫病、灰霉病、叶霉病 & 棉铃虫、白粉虱、斑潜蝇 \\
    柑橘 & 黄龙病、溃疡病、疮痂病、炭疽病 & 红蜘蛛、潜叶蛾、蚜虫 \\
    苹果 & 轮纹病、炭疽病、斑点落叶病、腐烂病 & 桃小食心虫、蚜虫、红蜘蛛 \\
    \bottomrule
  \end{tabular}
\end{table}

\section{附录2 市场调研核心数据}

\begin{table}[H]
  \centering
  \small
  \renewcommand{\arraystretch}{1.3}
  \caption{市场调研核心数据}
  \begin{tabular}{p{9.5cm}c}
    \toprule
    指标 & 数值 \\
    \midrule
    我国小农户数量（户） & 2 亿以上 \\
    农户智能手机使用率 & 76\% \\
    从未关注过农业类公众号的受访者占比 & 约 2/3 \\
    智慧农业市场规模（2024 年） & 约 8500 亿元 \\
    智慧农业市场规模（2028 年预测） & 约 14800 亿元 \\
    认为现有农业 APP 门槛过高、用不起来的受访者占比 & 73\% \\
    凭经验识别病虫害、存在识别不准的受访者占比 & 68\% \\
    常错过防治最佳窗口期的受访者占比 & 47\% \\
    \bottomrule
  \end{tabular}
\end{table}

\section{附录3 技术专利与成果说明}

\begin{enumerate}[leftmargin=1.6em]
  \item 发明专利（拟申请）：一种面向移动端的轻量化农作物病虫害图像识别方法及装置。
  \item 发明专利（拟申请）：一种基于熵权多属性决策的农作物农药精准推荐方法。
  \item 软件著作权（拟登记）：慧农识病移动端病虫害识别与防治应用软件。
  \item 软件著作权（拟登记）：基于联邦学习的分布式病虫害模型训练系统。
\end{enumerate}

\section{附录4 合作意向证明文件}

本项目已与以下机构建立初步合作意向（正式文件签署中，此处为占位说明）：

\begin{enumerate}[leftmargin=1.6em]
  \item 某县农业农村局——拟将慧农识病纳入当地数字乡村与统防统治示范项目；
  \item 某农资连锁集团——拟就农资渠道佣金分成与精准推广达成合作；
  \item 某农业大学植保学院——拟就病虫害防治知识库建设与模型评测达成产学研合作。
\end{enumerate}

\begin{center}
  \vspace{1em}
  \color{grayline}\rule{0.6\linewidth}{0.4pt}\\[0.5em]
  {\small （本商业计划书正文完）}
\end{center}
